\scanpage{36}
i.e.
\diagram{\begin{tikzcd}[ampersand replacement=\&,row sep=1.3em,column sep=0.8em]
D'{}^\circ \arrow[d,hook] \& \times \& \underline{\operatorname{Hom}}_{1/2\,\mathrm{add}}(\mathcal A,C') \arrow[r] \& C'\\
\underline{\operatorname{Hom}}_{1/2\,\mathrm{add}}(\mathcal A,\mathrm{Ab})^\circ \& \& \&
\end{tikzcd}}
noté
\[
 \underline{\operatorname{Hom}}_C(F,G)
\]

-\, -\, -\, -

\noindent\hfil\rule{0.4\linewidth}{0.35pt}\hfil

On suppose maintenant qu'on a
une catégorie $\mathscr S$, une sous-catégorie
\inserted{pl $C$,} de $\mathscr S_{\mathrm{ab}}$ stable par sommes directes
(alors $C$ est additive) et une
sous-catégorie pleine $\mathcal A$ de $C$ telle
que tt objet de $C$ soit isom à une
somme directe \inserted{dans $C$ donc dans $\mathscr S_{\mathrm{ab}}$} \deleted{\unsure{de}} d'objets de $\mathcal A$.

On a alors \deleted{\unsure{fonct.}}
\[
 \unsure{\alpha|u},\quad \mathscr S^\circ\longrightarrow
 \deleted{\widehat{\mathscr S}_{\mathrm{ab}}}
 \underline{\operatorname{Hom}}_{\mathrm{add}}(C,\mathrm{Ab})
\]
\deleted{\unsure{comme}} en associant à tt $X\in\operatorname{Ob}\mathscr S$
le foncteur $M\longmapsto\operatorname{Hom}_{\mathscr S}(X,M)$.
Or ce foncteur est connu à équiv.
près par \unsure{de --|sa --|quand on} \unsure{connaissant|en connaissant|connaît} son image
dans $\underline{\operatorname{Hom}}_{\unsure{.\mathrm{add}|\mathrm{add}|1/2\,\mathrm{add}}}(\mathcal A,(\mathrm{Ab}))$, donc l'étude
de $\alpha$ est \unsure{équivalente} à celle de
% archive sheet number: 35
